% Copyright 2026 Myeongjin Kang. Original prose CC BY 4.0; original code MIT.
\documentclass[11pt,a4paper]{article}
\usepackage[T1]{fontenc}
\usepackage{lmodern}
\usepackage[margin=1in]{geometry}
\usepackage{amsmath,amssymb,amsthm,mathtools,microtype,booktabs,array}
\usepackage{xcolor}
\usepackage[colorlinks=true,linkcolor=blue!50!black,citecolor=blue!50!black,urlcolor=blue!50!black]{hyperref}
\hypersetup{pdftitle={Doubling rectifiable curves with no finite-dimensional bi-Lipschitz embedding},pdfauthor={Myeongjin Kang},pdfsubject={Ambient doubling thickening in Hilbert and Banach spaces}}
\setlength{\emergencystretch}{2em}
\newtheorem{maintheorem}{Theorem}
\renewcommand{\themaintheorem}{\Alph{maintheorem}}
\newtheorem{theorem}{Theorem}[section]
\newtheorem{lemma}[theorem]{Lemma}
\newtheorem{proposition}[theorem]{Proposition}
\newtheorem{corollary}[theorem]{Corollary}
\newtheorem{premise}{Input theorem}
\theoremstyle{definition}
\newtheorem{definition}[theorem]{Definition}
\theoremstyle{remark}
\newtheorem{remark}[theorem]{Remark}
\numberwithin{equation}{section}
\newcommand{\R}{\mathbb R}
\newcommand{\N}{\mathbb N}
\newcommand{\norm}[1]{\lVert #1\rVert}
\newcommand{\diam}{\operatorname{diam}}
\newcommand{\HH}{\mathcal H}
\newcommand{\cB}{\overline B}
\title{Doubling rectifiable curves with no finite-dimensional bi-Lipschitz embedding}
\author{Myeongjin Kang}
\date{October 2026}
\begin{document}
\maketitle
\begin{abstract}
Using the fixed doubling Hilbert-space obstruction of OpenAI as an external premise, we construct in every infinite-dimensional real Banach space a Lipschitz image of $[0,1]$ that is uniformly doubling in its ambient norm metric and admits no bi-Lipschitz embedding into any finite-dimensional real normed space at any finite distortion. The same compact, path-connected set of finite one-dimensional Hausdorff measure works for all target dimensions and distortions. The additional construction is a dimension-independent ambient doubling thickening of a finite doubling Hilbert set: its dyadic nets are lifted to their own scales in one extra coordinate and joined by parent segments. Direct estimates control balls at edge interiors, including the bottom vertical edges and all scale regimes. Classical Dvoretzky embeddings of the whole finite graph spans, shrinking placement along a spine, and summable closed tours then produce the curve. No intrinsic path metric, injective parametrization, or regularity beyond the stated conclusions is asserted.
\end{abstract}

\section{Conventions, input, and main result}
All balls in this paper are \emph{closed}. If $(M,d)$ is a metric space, write
\[
 \cB_M(x,r)=\{y\in M:d(x,y)\le r\}.
\]
We call $M$ \emph{$D$-doubling} if every such ball with $x\in M$ and $r>0$ is covered by at most $D$ balls of radius $r/2$ with centers in $M$. Centers of a covering ball need not belong to the ball being covered. For a subset of a normed space the metric is always the restricted ambient norm, never the graph or path metric.

An embedding $f:M\to N$ has \emph{distortion at most $A$} if some scale $a>0$ satisfies
\begin{equation}\label{eq:distortion}
 a\,d_M(x,y)\le d_N(f(x),f(y))\le Aa\,d_M(x,y)
 \qquad(x,y\in M).
\end{equation}
The target dimension and the distortion are both allowed to be arbitrary finite numbers.

Our external nonembedding input is the following theorem, used as stated, without a proof here.
\begin{premise}[OpenAI, Theorem~1]\label{input:source}
There is a fixed subset $S$ of real $\ell_2$, with its induced Hilbert metric, that is $\lambda$-doubling for $\lambda=76800$ and admits no bi-Lipschitz embedding into $\R^k$ for any positive integer $k$, at any finite distortion.
\end{premise}
The source is the pinned September~25, 2026 preprint~\cite{OpenAI}. Its compact Banach-space corollary already supplies finite witnesses, Dvoretzky placement and a disconnected compact assembly. We reproduce the parts of those deductions needed here, with attribution, and use Input theorem~\ref{input:source} as the sole nonembedding premise.

Set
\begin{equation}\label{eq:constants}
 D=\lambda^2,\qquad \Delta=256D^{10}=256\lambda^{20},\qquad
 \Lambda=10+6\Delta^8.
\end{equation}
These are finite universal integers.
\begin{maintheorem}\label{thm:main}
Under Input theorem~\ref{input:source}, every infinite-dimensional real Banach space $B$ contains a set $\Gamma=\gamma([0,1])$, for a Lipschitz map $\gamma:[0,1]\to B$, such that:
\begin{enumerate}
\item $\Gamma$ is $\Lambda$-doubling in its restricted ambient norm metric;
\item $\Gamma$ admits no finite-distortion bi-Lipschitz embedding into any finite-dimensional real normed space.
\end{enumerate}
In particular, $\Gamma$ is one compact, path-connected set with $\HH^1(\Gamma)<\infty$, chosen before the target dimension and distortion are specified.
\end{maintheorem}
We do not claim that $\gamma$ is injective, that $\Gamma$ is a simple arc, that its path metric is doubling, or that it is uniformly quasiconvex or Ahlfors regular. The constants in~\eqref{eq:constants} are not optimized, and no quantitative cardinality bound for the finite witnesses is required.

The main additional ingredient is the finite ambient thickening of Section~\ref{sec:thickening}. Net hierarchies and scale separation are established tools, for example in Gupta--Talwar's intrinsic doubling completion~\cite{GT}. Here the net scale is made into an actual extra Hilbert coordinate. This controls ambient balls meeting long segment interiors, an issue not settled by doubling of a graph's shortest-path metric. Related distinctions are discussed in Section~\ref{sec:comparison}.

\section{Finite ambient doubling thickening}\label{sec:thickening}
For a real Hilbert space $H$, the sum $H\oplus\R$ carries its Hilbert product norm.
\begin{theorem}\label{thm:finite}
Let $X\subset H$ be finite, with at least two points, $\diam X=1$, and doubling constant at most an integer $D\ge1$. There is a finite abstract rooted tree whose straight-segment realization $G\subset H\oplus\R$ has the following properties:
\begin{enumerate}
\item $X\times\{0\}\subset G$, with all its original distances preserved;
\item $G$ is connected, lies in a ball of radius $\sqrt2$ about its root $o$, and has finite abstract edge-length sum;
\item $G$ is $256D^{10}$-doubling in its ambient Hilbert metric.
\end{enumerate}
The doubling constant is independent of $|X|$, its minimum separation and the dimension of its span.
\end{theorem}
We separate the construction from the ball estimates. Segment crossings or overlaps are permitted; the abstract tree is used for traversal and length accounting, not to define distances on $G$.

\subsection{The finite tree and its realization}
Choose $x_0\in X$, put $N_0=\{x_0\}$, and let $s_k=2^{-k}$. Choose $m\ge1$ such that $s_m$ is smaller than the minimum distance between distinct points of $X$. For $1\le k\le m$, let $N_k\subset X$ be maximal with distinct points at distances strictly greater than $s_k$. Maximality makes $N_k$ an $s_k$-net: every point of $X$ is within distance $s_k$ of $N_k$. Also $N_m=X$, and $N_0$ is a $1$-net since $\diam X=1$.

For each $x\in N_k$ choose a parent $p_k(x)\in N_{k-1}$ with
\begin{equation}\label{eq:parent}
 \norm{x-p_k(x)}\le s_{k-1}=2s_k.
\end{equation}
The abstract vertices are the labels $(k,x)$ for $0\le k\le m$, $x\in N_k$, together with distinct bottom labels $b_x$ for $x\in X$. Join $(k,x)$ to $(k-1,p_k(x))$ and $b_x$ to $(m,x)$. Each nonroot label has exactly one parent, and every parent chain reaches $(0,x_0)$. This finite connected graph is a rooted tree.

Realize $(k,x)$ as $(x,s_k)$ and $b_x$ as $(x,0)$, and realize each abstract edge as the straight segment between its endpoint images. Denote the union by $G$ and its root by $o=(x_0,1)$. The regular edge from scale $s=s_k$ has the parametrization
\begin{equation}\label{eq:edgeparam}
 \big((1-\theta)x+\theta p_k(x),(1+\theta)s\big),\qquad 0\le\theta\le1.
\end{equation}
Bottom edges are $\{(x,t):0\le t\le s_m\}$. Thus $G$ is connected and contains $X\times\{0\}$. Its projected points are convex combinations of points of $X$, so they are within distance one of $x_0$. All heights belong to $[0,1]$, proving $G\subset\cB_{H\oplus\R}(o,\sqrt2)$.

Let $L(G)$ be the sum of lengths of the realized \emph{abstract} edges, counted once per abstract edge, even if their images overlap. By~\eqref{eq:parent},
\begin{equation}\label{eq:length}
 L(G)\le \sqrt5\sum_{k=1}^{m}s_k|N_k|+s_m|X|<\infty.
\end{equation}
Indeed a regular edge has vertical displacement $s_k$ and horizontal displacement at most $2s_k$. We only use $\HH^1(G)\le L(G)$; equality is neither needed nor asserted.

\subsection{Packing and segment covers}
\begin{lemma}\label{lem:packing}
Let $R>0$. If $P\subset\cB_X(x_*,R)$ has pairwise distances strictly greater than $s>0$, then
\[
 |P|\le D^q\quad\hbox{whenever}\quad \frac{2R}{2^q}\le s.
\]
Here $x_*\in X$ and $q$ is a nonnegative integer.
\end{lemma}
\begin{proof}
Iterating the doubling cover $q$ times gives at most $D^q$ balls of radius $R/2^q$, each of diameter at most $s$. Each contains at most one point of $P$.
\end{proof}
If a closed ambient ball of radius $r$ meets a segment, their intersection is a closed subsegment, possibly one point, of length at most $2r$. It can be covered by two closed balls of radius $r/2$ centered on that subsegment, by placing centers at its quarter points. We retain the safe bound of eight such balls per edge in the counts below.

Fix a ball of radius $r>0$ about $(p,h)\in G$. No vertex assumption is made about this center. Every point of the ball's intersection with $G$ lies on a bottom edge or a regular edge, and we cover these portions separately. All cover centers will lie in the portion being covered, hence in $G$, except where an explicit larger-set cover is stated.

\subsection{Bottom vertical edges: every height}
The projections of bottom-edge points in the ball lie in $X\cap\cB_H(p,r)$. If this set is nonempty, choose one of its members $x_*$. These projections lie in $\cB_X(x_*,2r)$, which four halvings cover by at most $D^4$ $X$-balls of radius $r/8$. Their heights lie in
\[
 [h-r,h+r]\cap[0,s_m],
\]
an interval of length at most $2r$, covered by at most sixteen closed intervals of length at most $r/8$. Group bottom points according to a projection ball and a height interval. Within one group the squared distance is at most
\begin{equation}\label{eq:bottom}
 (r/4)^2+(r/8)^2=5r^2/64<r^2/4.
\end{equation}
Choose one point from each nonempty group as a center. The whole bottom portion is covered by at most $16D^4$ balls of radius $r/2$. This estimate does not depend on $h$ or on the last scale $s_m$.

\subsection{Low centers: fine and intermediate regular edges}
Suppose $h\le2r$. A meeting regular edge of scale $s$ has $s\le h+r\le3r$, by its height interval $[s,2s]$.

\paragraph{Fine edges: $s\le r/64$.}
By~\eqref{eq:edgeparam}, a point $q$ on such an edge is within horizontal distance $2s$ and height $2s$ of the bottom point $(x,0)$ associated with its lower endpoint. Consequently
\begin{equation}\label{eq:fineclose}
 \norm{q-(x,0)}\le\sqrt8s\le\frac{\sqrt8}{64}r<\frac r{16}.
\end{equation}
If the edge meets the ball, $\norm{x-p}<17r/16<2r$. Choose one such endpoint $x_*\in X$. All lower endpoints of meeting fine edges lie in $\cB_X(x_*,4r)$, which nine halvings cover by at most $D^9$ balls of radius $r/128$.

Assign each fine point in the original ball to one incident fine edge and then to a covering ball containing that edge's lower endpoint. Two points assigned to the same ball are at distance strictly less than
\begin{equation}\label{eq:finegroup}
 2(r/16)+2(r/128)=9r/64<r/2.
\end{equation}
A representative from each nonempty group provides a half-radius cover using at most $D^9$ centers. Endpoints from different fine scales are grouped together; no separation between those scales is assumed.

\paragraph{Intermediate edges: $r/64<s\le3r$.}
There are at most eight dyadic scales in this interval: nine distinct dyadic values would have largest-to-smallest ratio at least $2^8=256$, whereas the present ratio is strictly less than $192$.
At a fixed scale $s$, the lower endpoint of a meeting edge satisfies $\norm{x-p}\le2s+r$. Choosing one such endpoint $x_*\in N_k\subset X$ puts all of them in
\[
 \cB_X(x_*,4s+2r),\qquad 4s+2r<132s.
\]
Nine halvings give covering balls of diameter strictly less than $264s/512<s$. Since $N_k$ is strictly $s$-separated, Lemma~\ref{lem:packing} bounds the number of edges by $D^9$. Cover each edge portion with at most eight balls. Over all intermediate scales the count is at most $64D^9$.

Together with the bottom portion, every low-center ball needs at most
\begin{equation}\label{eq:lowcount}
 65D^9+16D^4\le81D^{10}
\end{equation}
half-radius balls. The boundaries $h=2r$ and $s=r/64$ are included in the stated cases.

\subsection{High centers}
Suppose $h>2r$. A meeting regular scale satisfies
\begin{equation}\label{eq:highwindow}
 \frac{h-r}{2}\le s\le h+r.
\end{equation}
The endpoint ratio of this interval is $2(h+r)/(h-r)<6$. In particular at most four dyadic scales occur; this is a safe overcount. Also $s>r/2$.
At a fixed scale, the lower endpoints satisfy $\norm{x-p}\le2s+r<4s$. Relative to one chosen endpoint $x_*\in X$, they lie in $\cB_X(x_*,8s)$. Five halvings have diameter $16s/32<s$, so at most $D^5$ separated endpoints, and hence edges, occur. Counting eight half-radius balls per edge gives at most $32D^5$ balls over the regular portion. Including the bottom portion yields
\begin{equation}\label{eq:highcount}
 32D^5+16D^4\le48D^{10}.
\end{equation}
Both~\eqref{eq:lowcount} and~\eqref{eq:highcount} are below $256D^{10}$, proving Theorem~\ref{thm:finite}. Every edge level, including the top parent endpoints and the bottom vertical edges, is included. Crossings and overlaps only create harmless overcounting. This proof concerns ambient balls and makes no assertion about the doubling constant of the tree's path metric.

\section{Finite witnesses and placement of the whole graph}
\subsection{Finite witnesses from the fixed source obstruction}
The following compactness deduction is the finite-witness argument in the proof of the source's compact Banach-space corollary~\cite{OpenAI}.
\begin{lemma}\label{lem:witness}
Under Input theorem~\ref{input:source}, for every positive integer $k$ and finite $A\ge1$ there is a finite subset of $S$ admitting no embedding into $\R^k$ with distortion at most $A$.
\end{lemma}
\begin{proof}
Suppose every finite subset admitted such an embedding. Fix $o\in S$ and let
\[
 Q_x=\{u\in\R^k:\norm u\le A\,d_S(o,x)\}\qquad(x\in S).
\]
Each $Q_x$ is compact, so $\prod_{x\in S}Q_x$ is compact. In this product impose the closed constraints
\[
 d_S(x,y)\le\norm{f(x)-f(y)}\le A\,d_S(x,y).
\]
Any finite list of constraints involves finitely many points. Embed those points together with $o$, translate to send $o$ to zero, and divide by the embedding's lower scale in~\eqref{eq:distortion}. The resulting coordinates belong to their $Q_x$; set all remaining coordinates to zero. This proves the finite-intersection property. Product compactness then produces an embedding of all of $S$ with distortion at most $A$, contradicting Input theorem~\ref{input:source}.
\end{proof}
Choose for each $j\ge1$ a finite $X_j\subset S$ with no embedding into $\R^j$ of distortion at most $j$. These choices are made once, before any final target is considered. Each $X_j$ has at least two points.

Every $X_j$ is $\lambda^2$-doubling in its induced metric. Indeed, an $X_j$-ball of radius $r$ lies in the corresponding $S$-ball. Two halvings cover the latter by $\lambda^2$ balls of radius $r/4$. In each ball meeting the original $X_j$-ball, choose a point of that intersection. Recenter there; the part in question is within distance $r/2$ of the new center. This keeps all centers in $X_j$.

Rescale each finite witness to diameter one, and retain the notation $X_j$. Rescaling preserves both doubling and the obstruction because distortion permits a free scale. Apply Theorem~\ref{thm:finite} with $D=\lambda^2$ to obtain $G_j$, root $o_j$, and abstract edge-length sum $L_j=L(G_j)$. All $G_j$ are $\Delta$-doubling, with $\Delta$ as in~\eqref{eq:constants}.

\subsection{Dvoretzky on the entire finite Euclidean span}
We use the classical finite-dimensional Dvoretzky theorem~\cite[Section~7]{Dvoretzky}: for every integer $n\ge1$, every infinite-dimensional real Banach space contains a linear copy of $\ell_2^n$ with distortion strictly smaller than two. A modern statement is recalled in~\cite[Theorem~1.1]{COO}. Only this qualitative finite-dimensional form is used.

For the given $B$, apply it to the finite-dimensional Hilbert space
\[
 E_j=\operatorname{span}(G_j-o_j),
\]
not just to the bottom endpoints. Since the graph has finitely many vertices and straight edges, $E_j$ is finite-dimensional. After a scalar normalization there is a linear map $T_j:E_j\to B$ with
\begin{equation}\label{eq:dvoretzky}
 \norm z\le\norm{T_jz}_B\le2\norm z\qquad(z\in E_j).
\end{equation}
Put $Y_j=T_j(G_j-o_j)$. It contains zero and is a connected finite union of segments, contained in $\cB_B(0,2\sqrt2)$. Its abstract edge-length sum is at most $2L_j$.

\begin{lemma}\label{lem:transfer}
Each $Y_j$ is $\Delta^4$-doubling in the ambient norm metric.
\end{lemma}
\begin{proof}
Let a $Y_j$-ball of radius $r$ be centered at $T_j(z-o_j)$, with $z\in G_j$. By the lower inequality in~\eqref{eq:dvoretzky}, its preimage lies in $\cB_{G_j}(z,r)$. Four halvings cover that ball by at most $\Delta^4$ $G_j$-balls of radius $r/16$. In each that meets the preimage choose a point of the intersection. Any two preimage points in the same covering ball are at Hilbert distance at most $r/8$, and their images are at distance at most $r/4<r/2$. The chosen images are valid centers in $Y_j$ for the desired cover.
\end{proof}

\section{A single rectifiable curve}
Fix a unit vector $v\in B$, put $a_j=2^{-j}$, and set $\eta=1/100$. Choose positive $\rho_j$ so small that
\begin{equation}\label{eq:scales}
 2\sqrt2\rho_j\le\eta a_j,
 \qquad 2\rho_jL_j\le2^{-j}.
\end{equation}
Define
\[
 C_j=a_jv+\rho_jY_j,\qquad I=\{tv:0\le t\le1\},
 \qquad \Gamma=I\cup\bigcup_{j\ge1}C_j.
\]
Each $C_j$ contains its root $a_jv$, lies in $\cB_B(a_jv,\eta a_j)$, and is $\Delta^4$-doubling. Write $\ell_j$ for the sum of its realized abstract edge lengths, counting every abstract edge once. Then $\ell_j\le2\rho_jL_j\le2^{-j}$. Intersections between realized edges do not affect this definition.

The set $\Gamma$ is compact. Indeed, a sequence in $\Gamma$ either has a subsequence in $I$ together with finitely many clusters, a compact set, or has a subsequence lying in clusters whose indices tend to infinity. In the second case its norms tend to zero, and zero belongs to $I$.

A depth-first traversal of the finite abstract tree realizing $C_j$ starts and ends at its root, traverses each abstract edge twice, and covers the whole realized graph. Parametrizing each traversal by its ambient norm length gives a closed $1$-Lipschitz tour of duration $2\ell_j$. Coincident realized vertices or crossings cause no difficulty: the traversal is defined on the abstract tree and then mapped into $B$.

Starting at $v$, move down $I$ at unit speed, interrupting the motion at each $a_jv$ to take this tour of $C_j$. The time immediately after the $j$th tour is
\[
 \tau_j=1-a_j+2\sum_{i=1}^j\ell_i,
 \qquad \tau_j\uparrow T:=1+2\sum_{i\ge1}\ell_i\le3.
\]
Every closed time interval contained in $[0,T)$ meets only finitely many pieces of this concatenation. Each finite concatenation is $1$-Lipschitz, so the resulting map is $1$-Lipschitz on $[0,T)$. After time $\tau_j$, its image lies in $\cB_B(0,(1+\eta)a_j)$: the remaining spine lies below $a_jv$, and every remaining cluster has smaller root coordinate. It therefore extends continuously to time $T$ with value zero. The Lipschitz inequality extends to this endpoint by passage to the limit.

Rescale time to $[0,1]$ and call the resulting map $\gamma$. Its image is exactly $\Gamma$, and $\operatorname{Lip}(\gamma)\le T\le3$. In particular $\Gamma$ is path connected and has finite $\mathcal H^1$ measure, by the usual Lipschitz-image bound for Hausdorff measure. No injectivity of $\gamma$ is asserted.

\section{Doubling of the assembled image}
We prove a bound independent of the sequence of graphs and of $B$. Fix $x\in\Gamma$ and $r>0$. The part of $I$ in $\cB_B(x,r)$ is an interval of length at most $2r$, because a norm ball is convex and distance along $I$ is the ordinary distance of its scalar coordinates. It can be covered by two closed balls of radius $r/2$ with centers on that interval. We retain the harmless upper bound of five such balls.

All clusters with $a_j\le r/4$ lie in the single ball $\cB_\Gamma(0,r/2)$, since
\[
 \norm z\le(1+\eta)a_j\le\frac{101}{400}r<\frac r2
 \qquad(z\in C_j).
\]
This center is permitted even if it lies outside the original ball.

It remains to count the clusters with $a_j>r/4$ that meet $\cB_B(x,r)$. If $i<j$, then $a_j\le a_i/2$, and any $z_i\in C_i$, $z_j\in C_j$ satisfy
\begin{align}
 \norm{z_i-z_j}
 &\ge a_i-a_j-\eta(a_i+a_j)\notag\\
 &\ge \frac{1-3\eta}{2}a_i=\frac{97}{200}a_i.
 \label{eq:cluster-separation}
\end{align}
If at least two of the remaining clusters meet the ball, take the one with largest root coordinate, say $a_i$. Points from it and from another intersecting cluster have distance at most $2r$, so~\eqref{eq:cluster-separation} implies $a_i\le400r/97$. Thus every intersecting root coordinate belongs to
\[
 (r/4,400r/97],
\]
whose endpoint ratio is $1600/97<32$. There are at most five dyadic numbers $2^{-j}$ in this interval: six such numbers would have largest-to-smallest ratio at least $32$. If there is at most one intersecting cluster, the same count of five is immediate without imposing the displayed upper bound on its root coordinate.

For each of these at most five clusters, choose $y_j\in C_j\cap\cB_B(x,r)$. Its intersection with the original ball lies in $\cB_{C_j}(y_j,2r)$. Two iterations of the $\Delta^4$-doubling property cover this by at most $\Delta^8$ balls of radius $r/2$, with centers in $C_j$. Altogether we have used at most
\[
 5+1+5\Delta^8=6+5\Delta^8\le10+6\Delta^8=\Lambda
\]
balls. This proves the claimed ambient doubling property at every radius and every center, including centers in the interiors of cluster edges.

\section{Excluding every finite-dimensional target}
The construction of $\Gamma$ used the entire sequence $(X_j)$ before any target space or distortion was chosen. Suppose that $f:\Gamma\to V$ is a finite-distortion bi-Lipschitz embedding into a finite-dimensional real normed space. The case $\dim V=0$ is impossible because $I$ is nontrivial. If $\dim V=k\ge1$, choose a linear isomorphism $V\to\mathbb R^k$. Composing with $f$ gives a bi-Lipschitz map of some finite distortion $A$ into Euclidean $\mathbb R^k$.

Choose an integer $j\ge k$ with $j\ge2A$. The bottom copy of $X_j$ in $G_j$ is isometric to $X_j$. Its image in $C_j$ under translation, scaling, and $T_j$ has distortion at most two by~\eqref{eq:dvoretzky}; the scalar $\rho_j$ does not affect distortion. Restrict the composite map to that copy and then include $\mathbb R^k$ isometrically in $\mathbb R^j$. This embeds $X_j$ into $\mathbb R^j$ with distortion at most $2A\le j$, contrary to its choice. This completes the proof of Theorem~\ref{thm:main}.

\section{Relation to earlier constructions}\label{sec:comparison}
The nonembedding input, extraction of finite witnesses, qualitative use of Dvoretzky's theorem, and separated accumulation of shrinking clusters are antecedents. In particular, the pinned source~\cite{OpenAI} already proves a compact doubling nonembedding subset of every infinite-dimensional Banach space. The additional operation here is Theorem~\ref{thm:finite}: a connected finite polygonal thickening of each finite Hilbert witness, retaining its original distances and a uniform doubling bound for the \emph{ambient Hilbert metric}. Summable tree tours then supply a single rectifiable curve.

Gupta and Talwar~\cite[Theorem~13 and Section~4]{GT} construct low-dimensional completions of doubling metrics using graph and shortest-path structures. Their net-tree and long-edge modifications address the metric of a completion. They do not directly give the ambient Hilbert realization needed here: mapping an abstract graph into Hilbert space can bring interiors of edges close together even when their endpoints are far apart in the graph metric. The height coordinate and the four ambient-ball cases above control precisely this issue.

Semmes~\cite[Section~6, especially Theorem~6.3]{Semmes} gives a closely related but different construction. A snowflake set in Euclidean space is used to define a weight involving distance to that set, and a weighted path metric produces a doubling metric space with curves. That metric is not the inherited constant-coefficient Hilbert metric. Consequently it does not directly provide a finite-dimensional normed span to which the Dvoretzky step in our argument applies. This comparison concerns the actual weighted-metric construction, not merely the common presence of curves and doubling.

Schul's Hilbert-space traveling-salesman theorem~\cite[Theorem~1.5 and Section~4]{Schul} supplies rectifiable connected supersets from a multiscale beta-number estimate. Its stated conclusion controls length; it does not assert a universal doubling bound for the ambient metric of the output set in terms of a doubling constant of the input. The estimate in Theorem~\ref{thm:finite} supplies that additional property directly. We make no claim of an optimal constant, an injective curve, or an exhaustive classification of earlier connected metric constructions. These comparisons are bounded comparisons of the indicated primary results and proof operations, not a certification of worldwide priority.

\section{Proof status and reproducibility}
The accompanying Lean package, \texttt{CDFLocal}, checks selected metric-cover interfaces and exact bounds. It defines closed-ball covers by actual finite sets of centers and proves cover composition, iterated doubling, recentering onto subsets, separated-set packing, $D^2$ doubling for metric subsets, and $D^4$ doubling under a surjective noncontracting $2$-Lipschitz map. It also checks dyadic scale counts, a five-center normed-segment cover, parent-edge projection and height inequalities, metric grouping and cluster-separation estimates, and the polynomial count bounds used here. The segment lemma is a sufficient five-center bound; the elementary two-center observation in this manuscript is not its formal statement.

These are conditional support results. The package does not construct the net tree, realize the complete Hilbert graph, prove the whole all-regime graph-cover theorem, instantiate Dvoretzky, construct the infinite tour, prove compactness, or derive the final nonembedding theorem. In particular, the geometric hypotheses connecting individual inequalities to the constructed graph are supplied by the written proof. The source theorem is also external. The exact formal source, declaration types, coverage limits, and validation records are separately hash-bound in the accompanying provenance packet. No full formalization of Theorem~\ref{thm:main} is claimed.

The mathematical proof in this article is a human-readable deduction from Input Theorem~\ref{input:source} and the cited classical results. The source theorem is an explicitly accepted external premise; its nonembedding proof is not re-proved or audited here. A separate source-and-proof map records the exact pinned source interfaces, the reviewed scope, and the limits of the literature reading. The accompanying build script reproduces this PDF from the original TeX source. Integrity checks authenticate file contents and review dependencies; they are not mathematical proof certificates for the full theorem.

\paragraph{Author and license.}
Myeongjin Kang; no institutional affiliation is claimed. The original prose is licensed under Creative Commons Attribution 4.0 International (CC BY 4.0). Original accompanying code is licensed under the MIT License. Cited third-party works retain their own rights and are not relicensed or redistributed as part of the publication artifact.

\begin{thebibliography}{9}
\bibitem{OpenAI}
OpenAI, \emph{A doubling Hilbert subset with no finite-dimensional bi-Lipschitz embedding}, September 25, 2026. Pinned manuscript in the OpenAI math repository, revision \texttt{fd4aeeb2ee4fc729c18d98444fed42fd0529eeeb}; main theorem and compact Banach-space corollary. \url{https://github.com/openai/math}.

\bibitem{Dvoretzky}
A. Dvoretzky, \emph{Some results on convex bodies and Banach spaces}, in \emph{Proceedings of the International Symposium on Linear Spaces (Jerusalem, 1960)}, Jerusalem Academic Press and Pergamon Press, 1961, pp.~123--160.

\bibitem{COO}
F. Catrina, S. Ostrovska, and M. I. Ostrovskii, \emph{Dvoretzky-type theorem for locally finite subsets of a Hilbert space}, Annales de l'Institut Fourier \textbf{75} (2025), no.~6, 2565--2607. \url{https://doi.org/10.5802/aif.3672}.

\bibitem{GT}
A. Gupta and K. Talwar, \emph{How to Complete a Doubling Metric}, arXiv:0712.3331v2, 2007. \url{https://arxiv.org/abs/0712.3331v2}.

\bibitem{Semmes}
S. Semmes, \emph{On the nonexistence of bilipschitz parameterizations and geometric problems about $A_\infty$-weights}, Revista Matem\'atica Iberoamericana \textbf{12} (1996), no.~2, 337--410. \url{https://doi.org/10.4171/RMI/201}.

\bibitem{Schul}
R. Schul, \emph{Subsets of Rectifiable Curves in Hilbert Space---The Analyst's TSP}, arXiv:math/0602675v2, 2007. \url{https://arxiv.org/abs/math/0602675v2}.
\end{thebibliography}
\end{document}
